\section{Background}
\label{sec:background}

\paragraph{Muon and its preconditioned variants.}
For a weight matrix $W\in\mathbb R^{m\times n}$ with gradient $G_t$, Muon
\citep{jordan2024muon} maintains momentum $M_t=\mu M_{t-1}+G_t$ with a Nesterov
source $U_t=G_t+\mu M_t$, and steps along an approximately orthogonalized
version of $U_t$. The orthogonalization is a fixed-degree polynomial iteration
rather than an exact polar factor: with $X_0=U_t/(\|U_t\|_F+\varepsilon)$ and
coefficients $(a,b,c)=(3.4445,-4.7750,2.0315)$,
\begin{equation}
\label{eq:ns}
X_{k+1}=aX_k+\left(bA_k+cA_k^2\right)X_k,\qquad A_k=X_kX_k^{\Tr},
\end{equation}
run for a fixed number of steps, ten in the recipe we pin. We write $\Nten$ for
this map and stress throughout that it is a finite, rounded iteration: our
measurements in Appendix~\ref{app:numerics} show its ten-step output sits at
cosine $0.979$ to a partial polar factor, so identities that hold for the exact
polar decomposition cannot be substituted for it.

Preconditioned variants whiten before orthogonalizing. Let
$C_t$ be an exponential moving average of the layer's input second moment,
$x x^{\Tr}$, sampled over token positions. For a PSD moment
$F\in\mathbb R^{d\times d}$ with positive trace, define the trace-normalized,
damped moment and its two roots
\begin{equation}
\label{eq:factors}
A_d(F)=\frac{dF}{\tr F}+\epsilon I_d,\qquad
\In_d(F)=A_d(F)^{-1/2},\qquad
\Out_d(F)=\frac{\sqrt d\,A_d(F)^{-1/8}}{\|A_d(F)^{-1/8}\|_F},
\end{equation}
with $\epsilon=10^{-3}$ and spectral (not elementwise) powers. The inner root
is unnormalized; the outer root is scaled to unit RMS. Writing
$R=\In_n(C_t)$ and $\hat R=\Out_n(C_t)$, the input-side baseline's update
is\footnote{\todo{citation for the pinned preconditioned-MuonT implementation
these runs adapt.}}
\begin{equation}
\label{eq:muont}
\Delta W_t=-\eta_t\Bigl[\,s\,\Nten\!\left(U_tR\right)\hat R
\;+\;\lambda\,W_t\hat R^{2}\Bigr],
\qquad s=\sqrt{\max(1,m/n)} .
\end{equation}
Both factors act on the \emph{right}, because both are built from input
statistics, and the decay term is shaped by the outer factor. The split powers
$(1/2,1/8)$ rather than $(1/4,1/4)$ are the pinned recipe's choice. This is the
optimizer we call input-side MuonT and treat as the baseline to beat.

\paragraph{The unused side.}
Since $G=\delta x^{\Tr}$, the layer's gradient statistics have two natural
second moments, and \eqref{eq:muont} uses one. Let $\delta$ be the cotangent at
the layer's output and
\begin{equation}
\label{eq:bmoment}
B_t=\beta B_{t-1}+(1-\beta)\frac{1}{N_t}\sum_{j=1}^{N_t}\delta_{t,j}\delta_{t,j}^{\Tr},
\end{equation}
sampled at the same token positions as $C_t$. Then $L=\In_m(B_t)$ and
$\hat L=\Out_m(B_t)$ are output-side analogues of $R$ and $\hat R$, and they act
on the \emph{left}. The whole family of updates we study is obtained by
choosing which of the four factors to apply:
\begin{equation}
\label{eq:family}
\Delta W_t=-\eta_t\Bigl[\,s\,\hat L\,\Nten\!\left(L\,U_t\,R\right)\hat R
\;+\;\lambda\,W_t\hat R^{2}\Bigr],
\end{equation}
with absent factors read as the identity. Equation~\eqref{eq:muont} is the case
$L=\hat L=I$; the arms of Section~\ref{sec:method} are the other cases.

\paragraph{What $B$ is and is not.}
It is worth being precise, because the motivation for \eqref{eq:family} is
easy to overstate and our results do not test the strong version of it. $B_t$ as
defined in \eqref{eq:bmoment} is an \emph{uncentered second moment of observed
cotangents}. It is not an expectation under the model's own predictive
distribution, so it is not automatically a Fisher or
generalized Gauss--Newton factor; it is the empirical-Fisher-style quantity
whose limitations are well documented \citep{kunstner2019limitations}. Treating
$B\otimes C$ as a curvature model additionally assumes a Kronecker
factorization, as K-FAC does \citep{martens2015kfac}, and that approximation
discards cross-token terms which may matter at these sequence lengths. We have
not measured whether it holds here --- Section~\ref{sec:program} lists the
check --- so throughout this paper we treat $B$ as a \emph{normalizing
statistic that happens to be available}, and our claims are about what
applying it does to the loss, not about why.

One further caveat shapes the arm design. An exact polar factor of a
full-row-rank wide matrix $P$ satisfies $PP^{\Tr}=I$, which would make a left
factor applied after orthogonalization look redundant. It is not redundant
here, for two reasons: $\Nten$ is finite and rounded rather than exact, and the
outer input factor $\hat R$ is applied after it, so the final direction's
geometry is not the polar factor's. This is also why we do not assume the
inner-left and outer-left factors play symmetric roles, and why arms that apply
only $\hat L$ are measured separately from arms that apply both.

\paragraph{Related constructions.}
Full-matrix and Kronecker-factored preconditioners
\citep{martens2015kfac,gupta2018shampoo,anil2020scalable} use both sides of the
gradient but do not orthogonalize, and the Shampoo family routinely accepts
preconditioner staleness of hundreds of steps; our freshness result
(Section~\ref{sec:freshness}) is a measurement of how much a four-update
staleness costs in the Muon setting, where the orthogonalization already
supplies much of the conditioning. SOAP \citep{vyas2024soap} and other recent
hybrids sit between these regimes. We discuss the relationship in
Section~\ref{sec:related}; the short version is that our arms are not
reproductions of any published method, and where a published form is close we
label our version a formula control rather than a reproduction.
